% !TeX spellcheck = ru_RU
% !TEX root = vkr.tex

\section{Описание реализации}
\label{sec:implementation}

В предыдущем разделе виртуальная машина рассматривалась как набор
компонентов и поток управления между ними. В данном разделе описаны детали
реализации, которые определяют поведение этих компонентов.

\subsection{Прототип виртуальной машины}

Разработка решения началась с прототипа исполнителя для существующего формата
байткода Lama. Цель этого этапа состояла в том, чтобы проверить возможность
непосредственного выполнения байткода, порождаемого существующим компилятором,
и определить минимальный набор механизмов, необходимых для его интерпретации.
Прототип поддерживал базовое стековое исполнение инструкций, передачу
управления, вызовы функций и работу с основными значениями языка через
существующую среду выполнения. Это позволило подтвердить практическую
исполнимость байткодного пути и выделить ключевые требования к состоянию
виртуальной машины.

Одновременно прототип показал ограничения исходного формата байткода. Для
самостоятельного запуска многомодульных программ в нем отсутствовали сведения,
необходимые для загрузки зависимостей, регистрации публичных символов и
разрешения внешних ссылок. Кроме того, исполнение исходного байткода без
предварительной подготовки оставляло проверки корректности и часть разрешения
ссылок на этапе выполнения. Поэтому следующим шагом стало расширение формата
байткода и фиксация его структуры в явной спецификации, после чего была
построена итоговая реализация виртуальной машины, в которой фазы загрузки,
связывания, подготовки программы и исполнения разделены явно.

\subsection{Доработка генерации байткода}

Виртуальная машина исполняет не исходное стековое представление компилятора, а
байткодный файл. Поэтому часть работы была связана с тем, чтобы сделать этот
файл пригодным для самостоятельной загрузки и связывания. В существующей
инфраструктуре байткод фактически был сериализацией стековых инструкций одного
модуля. Для виртуальной машины этого недостаточно: при запуске необходимо
знать, какие модули импортируются, какие имена экспортируются наружу, сколько
глобальных переменных требуется модулю и какие ссылки должны быть разрешены вне
текущего файла.

Итоговая структура байткодного файла показана на
рисунке~\ref{fig:bytecode-format}. Все многобайтовые целые числа, как и раньше,
записываются в порядке от младшего байта к старшем. Изменения касаются прежде
всего заголовка, служебных секций и кодирования отдельных инструкций.

\begin{figure}[H]
    \centering
    \resizebox{\textwidth}{!}{\begin{tikzpicture}[
    section/.style={
        draw,
        semithick,
        align=center,
        minimum height=14mm,
        minimum width=28mm,
        text width=24mm,
        outer sep=0pt,
        font=\small
    },
    header/.style={section, fill=gray!12},
    block/.style={section, fill=gray!12},
    codeblock/.style={section, fill=gray!12},
    field/.style={
        draw,
        semithick,
        align=center,
        minimum height=12mm,
        minimum width=32mm,
        text width=29mm,
        outer sep=0pt,
        fill=gray!12,
        font=\scriptsize
    },
    detail/.style={
        draw,
        semithick,
        align=center,
        minimum height=9mm,
        minimum width=34mm,
        text width=31mm,
        outer sep=0pt,
        fill=gray!12,
        font=\scriptsize
    },
    widefield/.style={
        draw,
        semithick,
        align=center,
        minimum height=12mm,
        minimum width=45mm,
        text width=42mm,
        outer sep=0pt,
        fill=gray!12,
        font=\scriptsize
    },
    codefield/.style={widefield, fill=green!17},
    note/.style={
        draw,
        semithick,
        align=center,
        rounded corners=2pt,
        minimum height=9mm,
        minimum width=64mm,
        text width=60mm,
        fill=gray!12,
        font=\scriptsize
    },
    attach/.style={->, semithick},
    weakattach/.style={->, semithick}
]

\path coordinate (mainaux) at (0,0)
    node[
        header,
        anchor=west,
        label={[name=main-l1, align=center, font=\small, text width=31mm]center:
            {заголовок\\24 байта}}
    ] (header) at (mainaux) {}
    node[
        block,
        anchor=west,
        label={[name=main-l2, align=center, font=\footnotesize, text width=31mm]center:
            {таблица строк\\переменный размер}}
    ] (strings) at (header.east) {}
    node[
        section,
        fill=green!17,
        anchor=west,
        label={[name=main-l3, align=center, font=\small, text width=31mm]center:
            {импорты\\$n_i \times 4$ байта}}
    ] (imports) at (strings.east) {}
    node[
        section,
        fill=green!17,
        anchor=west,
        label={[name=main-l4, align=center, font=\footnotesize, text width=31mm]center:
            {публичные \\ символы\\$n_p \times 9$ байт}}
    ] (publics) at (imports.east) {}
    node[
        codeblock,
        anchor=west,
        minimum width = 12mm,
        label={[name=main-l5, align=center, font=\small, text width=16mm]center:
            {байткод \\ до \texttt{0xFF}}}
    ] (code) at (publics.east) {};

\node[detail, fill=green!17] (h0) at ($(header.west)+(-2.8,2.25)$)
{магическое число: \\ \texttt{uint8[4]}};
\node[detail, fill=green!17, below=0pt of h0] (h1)
{версия: \texttt{int32}};
\node[detail, below=0pt of h1] (h2)
{размер таблицы \\ строк: \texttt{int32}};
\node[detail, below=0pt of h2] (h3)
{число глобальных \\ переменных: \texttt{int32}};
\node[detail, fill=green!17, below=0pt of h3] (h4)
{число импортов: \texttt{int32}};
\node[detail, below=0pt of h4] (h5)
{число публичных\\символов: \texttt{int32}};
\draw[attach] (header.west) -- (h3.east);

\path coordinate (impaux) at ($(imports.south)+(-3.2,-1.65)$)
    node[
        field,
        fill=green!17,
        anchor=west,
        label={[name=imp-label, align=center, font=\scriptsize, text width=42mm]center:
            {смещение имени юнита\\в таблице строк\\\texttt{int32}}}
    ] (impentry) at (impaux) {};
\draw[attach] (imports.south) to[out=-90,in=90] (impentry.north);

\path coordinate (pubaux) at ($(publics.north)+(-6.8,1.8)$)
    node[
        field,
        anchor=west,
        label={[name=pub-l1, align=center, font=\scriptsize, text width=29mm]center:
            {смещение имени\\ в таблице строк \\\texttt{int32}}}
    ] (pub1) at (pubaux) {}
    node[
        field,
        anchor=west,
        label={[name=pub-l2, align=center, font=\scriptsize, text width=29mm]center:
            {адрес функции\\или индекс\\\texttt{int32}}}
    ] (pub2) at (pub1.east) {}
    node[
        field,
        fill=green!17,
        anchor=west,
        label={[name=pub-l3, align=center, font=\scriptsize, text width=29mm]center:
            {флаг\\\texttt{uint8}}}
    ] (pub3) at (pub2.east) {};
\draw[attach] (publics.north) to[out=90,in=-90] (pub2.south);

\end{tikzpicture}
}
    \caption{Формат байткодного модуля Lama. Зеленым отмечены поля и секции, добавленные или измененные по сравнению с исходным форматом.}
    \label{fig:bytecode-format}
\end{figure}

Первая группа изменений относится к общему устройству файла. Появление
спецификации позволило явно зафиксировать свойства байткода, которые раньше
были неявно заданы только компилятором. В заголовок были
добавлены магическое число и версия формата байткода, чтобы загрузчик мог
распознавать корректный вход и явно отличать версии байткода в будущем.
Кроме того, в заголовке и теле файла появились данные, необходимые для
модульной загрузки: число импортов, число публичных символов и соответствующие
секции \verb=imports= и \verb=publics=. Секция \verb=imports= хранит имена
импортируемых модулей через смещения в таблице строк. Секция \verb=publics=
содержит экспортируемые сущности модуля; по сравнению с исходным форматом в нее
был добавлен флаг типа символа, различающий публичную функцию и глобальное
значение. Это различие нужно при связывании, поскольку адрес исполняемого кода
и индекс глобальной переменной обрабатываются по-разному.

Вторая группа изменений связана с самими инструкциями. Для ссылок на внешние
сущности было введено единое кодирование отрицательными операндами через
таблицу строк. Инструкции \verb=CALL= и \verb=CLOSURE= используют такие
операнды для внешних функций, а \verb=LD_GLO= и \verb=ST_GLO= --- для внешних
глобальных значений. Неотрицательные значения по-прежнему обозначают локальные
ссылки внутри текущего модуля.

Кроме того, были изменены отдельные инструкции исходного байткода. В \verb=BEGIN_CLOSURE= был
добавлен операнд, задающий число замыкаемых переменных и необходимый для валидации стека. Была добавлена инструкция \verb=FAIL_KEEP= для сохранения семантики выполнения программы. Наконец, несколько операций исходного байткода,
таких как \verb=read=, \verb=write=, \verb=length= и \verb=string=, в новой реализации отображаются
на прямые вызовы функций существующей среды выполнения.

Еще одна доработка связана с именованием. Префиксы, используемые нативными
бэкендами для ассемблерных символов функций, встроенных функций и глобальных
значений, были оставлены на стороне генерации нативного кода.
Стековое представление и байткод при этом работают с исходными именами
сущностей. Благодаря этому байткодный формат не зависит от соглашений
нативного генератора, а нативная сборка сохраняет свои прежние правила
именования. 

\subsection{Загрузка и связывание модулей}
Входом виртуальной машины является либо имя байткодного модуля, либо путь к
файлу с расширением \verb=.bc=. Во втором случае имя основного модуля
получается из имени файла, а директория этого файла автоматически добавляется в
пути поиска. Дополнительные директории задаются параметром \verb=-I=.

Загрузка модулей строится вокруг импортов, записанных в байткоде. Для каждого
модуля считывается список зависимостей, после чего загрузчик ищет
соответствующие файлы в заданных директориях. При этом загрузчик поддерживает стек текущей рекурсии: если модуль
пытается импортировать один из модулей, находящихся в процессе загрузки, такая
ситуация распознается как циклическая зависимость и приводит к
ошибке.

Результатом этой фазы является упорядоченный набор байткодных модулей. В него
сначала попадают зависимости, а затем модуль, который от них зависит. Такой
порядок упрощает дальнейшую подготовку программы: к моменту связывания
основного модуля публичные символы импортированных модулей уже могут быть
зарегистрированы. 

Связывание исполняемых целей выполняется через релокации в \verb=converter.c=.
На этапе декодирования цель инструкции еще может быть известна только как
смещение внутри текущего байткодного моудля, имя публичной функции другого
модуля или имя внешней функции среды выполнения. Поэтому конвертер сначала
записывает \enquote{заглушку} и добавляет запись релокации: индекс исправляемого
слота, имя цели при необходимости и вид цели. В реализации различаются три
вида таких записей: внутренняя цель текущего модуля, публичная функция другого
модуля и \enquote{заглушка} вызова внешней функции.

Откладывание этого шага связано с устройством внутреннего формата программы.
Каждый модуль сначала декодируется в собственный временный массив инструкций.
Только после этого все такие массивы копируются в единый \verb=all_code=, где
адреса инструкций становятся стабильными. На финальном этапе релокации
заменяют \enquote{заглушки} настоящими указателями: внутренние цели
интерпретируются относительно начала своего модуля в \verb=all_code=,
межмодульные вызовы берут уже вычисленный индекс из таблицы символов, а
внешние вызовы направляются на сгенерированную инструкцию интерфейса
внешних функций.

Отдельно обрабатываются глобальные значения. Для инструкций загрузки и записи
глобального значения конвертер обычно сразу вычисляет указатель на
соответствующую область значений, поэтому общий механизм релокаций используется
прежде всего для исполняемых целей: переходов, вызовов, замыканий и
\enquote{заглушек} интерфейса внешних функций.

\subsection{Конвертация байткода во внутреннее представление}

После загрузки байткодного файла его кодовая секция преобразуется из последовательности байтов во
внутренний массив инструкций и их операндов. Эта фаза отделяет внешний бинарный формат от
представления, с которым работает исполнитель: код операции заменяется ссылкой на обработчик,
операнды декодируются один раз и сохраняются в той же форме, что и инструкции.

Для каждой инструкции конвертер , считывает операнды и добавляет во временный массив один или
несколько элементов внутреннего представления. Например, инструкция с числовым операндом
преобразуется в элемент с обработчиком операции и следующий за ним элемент с уже разобранным числом
в том же формате. Инструкции перехода и вызова требуют дополнительной обработки. Для внутренних переходов конвертер
проверяет корректность целевого смещения в пределах текущего модуля. Если же окончательный адрес
цели зависит от размещения кода в общей программе, конвертер сохраняет место, которое позднее
уточняется на этапе релокации.

\begin{figure}[H] \centering \resizebox{\textwidth}{!}{\begin{tikzpicture}[
    x=1cm,
    y=1cm,
    source/.style={
        draw,
        line width=0.35pt,
        draw=black!30,
        minimum width=37mm,
        minimum height=24mm,
        fill=white
    },
    sourceback/.style={
        draw,
        line width=0.35pt,
        draw=black!25,
        minimum width=37mm,
        minimum height=24mm,
        fill=gray!6
    },
    linkbox/.style={
        draw,
        rounded corners=2pt,
        semithick,
        align=center,
        text width=28mm,
        minimum height=11mm,
        inner sep=3pt,
        fill=gray!10,
        font=\scriptsize
    },
    insn/.style={
        draw,
        semithick,
        align=left,
        text width=16mm,
        minimum height=4.2mm,
        inner sep=1pt,
        fill=orange!18,
        font=\fontsize{5.4}{5.9}\selectfont\ttfamily
    },
    operand/.style={
        draw,
        semithick,
        align=left,
        text width=16mm,
        minimum height=4.2mm,
        inner sep=1pt,
        fill=orange!6,
        font=\fontsize{5.4}{5.9}\selectfont\ttfamily
    },
    operandwide/.style={
        operand
    },
    marker/.style={
        draw,
        semithick,
        align=center,
        text width=16mm,
        minimum height=4.2mm,
        inner sep=1pt,
        fill=gray!8,
        font=\fontsize{5.4}{5.9}\selectfont\ttfamily
    },
    flow/.style={->, semithick}
]

\node[sourceback] at (-4.23,0.44) {};
\node[sourceback] at (-4.34,0.32) {};
\node[source] (bc) at (-4.45,0.2) {};
\node[anchor=north west, inner sep=0pt] at ($(bc.north west)+(0.8mm,-0.8mm)$) {
    \begin{tikzpicture}[x=1mm,y=-1mm]
        \clip (0,0) rectangle (34.4,22.4);
        \node[
            anchor=north west,
            align=left,
            font=\ttfamily\fontsize{5.0}{5.6}\selectfont
        ] at (0.1,0.25)
            {PUBLIC ("main")\\
             LABEL ("main")\\
             BEGIN ("main", 2, 2, [], [], [])\\
             LD (Global ("x"))\\
             CONST (3)\\
             CALL ("func", 0, false)\\
             ST (Global ("z"))\\
             BINOP ("+")\\
             DROP\\
             ...};
    \end{tikzpicture}
};

\coordinate (convStart) at ($(bc.east)+(0.25,0)$);
\coordinate (convEnd) at (-0.20,0.2);
\draw[flow] (convStart) -- (convEnd)
    node[midway, above, font=\small] {конвертер};

\node[linkbox] (links) at (-1.15,-2.35)
    {опережающие ссылки\\
     соответствие индексов};

\node[insn] (c1) at (1.45,1.30) {\texttt{op\_begin}};
\node[operand, anchor=west] (a1b) at (c1.east) {\texttt{2}};
\node[operand, anchor=west] (a1c) at (a1b.east) {\texttt{2}};

\node[insn, below=0pt of c1] (c2) {\texttt{op\_ld\_glo}};
\node[operand, anchor=west] (a2) at (c2.east) {\texttt{\&globals[0]}};

\node[insn, below=0pt of c2] (c3) {\texttt{op\_const}};
\node[operand, anchor=west] (a3) at (c3.east) {\texttt{3}};

\node[insn, below=0pt of c3] (c4) {\texttt{op\_call}};
\node[operand, anchor=west] (a4a) at (c4.east) {\texttt{func}};
\node[operand, anchor=west] (a4b) at (a4a.east) {\texttt{0}};

\node[insn, below=0pt of c4] (c5) {\texttt{op\_st\_glo}};
\node[operand, anchor=west] (a5) at (c5.east) {\texttt{\&globals[2]}};

\node[insn, below=0pt of c5] (c6) {\texttt{op\_add}};

\node[insn, below=0pt of c6] (c7) {\texttt{op\_drop}};

\node[marker, anchor=north west] (dots) at (c7.south west) {\texttt{...}};

\draw[flow] (convEnd) -- (links.north);
\draw[flow] (convEnd) -- (c3.west);

\end{tikzpicture}
}
\caption{Преобразование кодовой секции байткодного модуля во внутреннее представление инструкций.}
\label{fig:bytecode-conversion} \end{figure}

% TODO: next three paragraphs haven't read them
На рисунке~\ref{fig:bytecode-conversion} показан процесс конвертации байткодных модулей.
Конвертер последовательно обходит кодовую секцию каждого загруженного юнита, формирует временный
массив инструкций и записывает релокации для тех операндов, которые нельзя заменить окончательным
адресом сразу. Запись релокации содержит индекс исправляемого элемента внутреннего массива, имя цели
при необходимости и вид ссылки: внутренняя цель текущего модуля, публичная функция другого модуля или
внешняя функция среды выполнения.

На этом же этапе выполняются проверки, которые зависят от структуры кодовой секции. Конвертер
проверяет, что цели переходов указывают на допустимые границы инструкций.
Для инструкций вызова и создания замыкания отдельно проверяется, что цель соответствует определенной в спецификации инструкции. Такие ошибки можно обнаружить непосредственно до выполнения самой программы.

Результатом конвертации является временное внутреннее представление отдельного модуля. После
обработки всех модули эти массивы копируются в общий массив инструкций виртуальной машины, где
адреса становятся стабильными и могут быть использованы при окончательном разрешении переходов,
вызовов и внешних ссылок.


\subsubsection{Стек виртуальной машины}

Стек виртуальной машины выделяется отдельно от системного стека. В текущей
реализации \verb=vm.c= создает область фиксированного размера с помощью
\verb=mmap=, использует верх этой области как базу стека, размещает глобальные
переменные непосредственно перед ней и начинает исполнение с вершиной стека,
указывающей на область глобальных значений. Эта же граница передается
существующему сборщику мусора, чтобы он просматривал стек виртуальной машины.

Изначально рассматривался вариант с выделением стека через
\verb=alloca=. На практике он оказался ненадежным при взаимодействии со средой
выполнения Lama: в области системного стека оставались устаревшие значения,
которые сборщик мусора мог принять за корни. Из-за этого в множество
достижимых значений попадали данные, уже не являвшиеся актуальными значениями
программы, что приводило к падениям. Отдельная область, выделенная через
\verb=mmap=, делает границы сканирования явными и ограничивает их стеком
виртуальной машины и глобальными переменными, предоставляя при этом динамическое расширение по мере
исполнения программы.

\subsubsection{Внутренний шитый код и внешние вызовы}

Внутренний массив инструкций организован так, чтобы каждый элемент содержал
либо ссылку на обработчик операции, либо ее уже подготовленный операнд. После
выполнения очередной операции управление переходит к следующему элементу или к
заранее вычисленной цели перехода. Такой подход уменьшает работу, выполняемую
в интерпретации: исполнитель не декодирует байткод повторно, а
работает с подготовленным представлением программы.

Схематично такой способ исполнения показан на
рисунке~\ref{fig:internal-threaded-code}. Подготовленное представление состоит
из основного массива \verb=all_code= и отдельного массива стартовых вызовов
\verb=entry_points=. Для каждого юнита в этом массиве хранится небольшой шаг
запуска: вызов его основной функции, число аргументов, удаление результата и
переход к следующему юниту. Первый такой шаг указывает на начало основного
кода первого юнита. Внутри \verb=all_code= массив \verb=insn= чередует элементы с
указателями на обработчики и элементы с операндами; обработчик сам знает, сколько
операндов необходимо прочитать.

Переход между обработчиками инструкций реализован как хвостовой вызов
следующего обработчика. Для этого в макросах диспетчеризации используется
атрибут компилятора \verb=__attribute__((musttail))=: он требует, чтобы вызов
следующего обработчика был скомпилирован как хвостовой переход. Без такого
требования последовательное исполнение инструкций могло бы накапливать кадры
на системном стеке и приводить к его переполнению на длинных программах. Кроме
того, хвостовой переход уменьшает накладные расходы диспетчеризации, так как
между обработчиками не выполняется обычная цепочка вложенных вызовов и
возвратов.

\begin{figure}[H]
    \centering
    \resizebox{\textwidth}{!}{\begin{tikzpicture}[
    x=1cm,
    y=1cm,
    entrycell/.style={
        draw,
        semithick,
        align=left,
        text width=31mm,
        minimum height=5.8mm,
        inner sep=2pt,
        font=\scriptsize
    },
    codecell/.style={
        draw,
        semithick,
        align=left,
        text width=39mm,
        minimum height=5.8mm,
        inner sep=2pt,
        font=\scriptsize
    },
    entryfunc/.style={entrycell, fill=orange!18},
    entryoperand/.style={entrycell, fill=orange!6},
    codefunc/.style={codecell, fill=orange!18},
    codeoperand/.style={codecell, fill=orange!6},
    markerentry/.style={entrycell, fill=gray!8, align=center},
    markercode/.style={codecell, fill=gray!8, align=center},
    handler/.style={
        draw,
        rounded corners=2pt,
        semithick,
        align=left,
        text width=47mm,
        minimum height=8mm,
        inner sep=4pt,
        fill=white,
        font=\scriptsize
    },
    ipbox/.style={
        draw,
        rounded corners=2pt,
        semithick,
        align=center,
        minimum width=8mm,
        minimum height=5mm,
        inner sep=4pt,
        fill=gray!8,
        font=\scriptsize
    },
    flow/.style={->, semithick},
    targetflow/.style={->, semithick, rounded corners=4pt},
    returnflow/.style={->, semithick, densely dashed}
]

\node[font=\small\bfseries, align=center] at (-4.3,0.85) {\\\texttt{entry\_points}};
\node[font=\small\bfseries, align=center] at (0.75,0.85) {\\\texttt{all\_code}};
\node[font=\small\bfseries] at (8.1,0.65) {\texttt{ops.c}};

\node[entryfunc] (entrycall1) at (-4.3,0) {\texttt{func: op\_call}};
\node[entryoperand, below=0pt of entrycall1] (entrytarget1) {\texttt{target: main\_1}};
\node[entryoperand, below=0pt of entrytarget1] (entryargs1) {\texttt{num: 0}};
\node[entryfunc,    below=0pt of entryargs1] (entrydrop1) {\texttt{func: op\_drop}};
\node[markerentry,  below=0pt of entrydrop1] (entrygap) {$\cdots$};
\node[entryfunc,    below=0pt of entrygap]   (entrycalln) {\texttt{func: op\_call}};
\node[entryoperand, below=0pt of entrycalln] (entrytargetn) {\texttt{target: main\_n}};
\node[entryoperand, below=0pt of entrytargetn] (entryargsn) {\texttt{num: 0}};
\node[entryfunc,    below=0pt of entryargsn] (entrydropn) {\texttt{func: op\_drop}};
\node[entryfunc,    below=0pt of entrydropn] (entryeof) {\texttt{func: op\_eof}};


\node[codefunc] (mainbegin) at (0.75,0) {\texttt{func: op\_begin (main\_1)}};
\node[codefunc,    below=0pt of mainbegin]  (constop) {\texttt{func: op\_const}};
\node[codeoperand, below=0pt of constop]    (constarg) {\texttt{num: 42}};
\node[codefunc,    below=0pt of constarg]   (callop) {\texttt{func: op\_call}};
\node[codeoperand, below=0pt of callop]     (calltarget) {\texttt{target: f}};
\node[codeoperand, below=0pt of calltarget] (callargs) {\texttt{num: n}};
\node[codefunc,    below=0pt of callargs]   (aftercall) {следующий обработчик};
\node[markercode,  below=0pt of aftercall]  (gap2) {$\cdots$};
\node[codefunc,    below=0pt of gap2]       (fbegin) {\texttt{func: op\_begin (f)}};
\node[markercode,  below=0pt of fbegin]     (fbody) {$\cdots$ тело функции};
\node[codefunc,    below=0pt of fbody]      (endop) {\texttt{func: op\_end}};

\node[ipbox] (ip) at ($(callop.west)+(-0.9,0)$) {\texttt{ip}};

\node[handler] (hcall) at (8.1,-1.25)
    {\texttt{void op\_call(DECL\_STATE) \{}\\
     \texttt{\hspace*{1.2em}ip++;}\\
     \texttt{\hspace*{1.2em}target = ip->target; ip++;}\\
     \texttt{\hspace*{1.2em}n\_args = ip->num; ip++;}\\
     \texttt{\hspace*{1.2em}PUSH\_FRAME(..., ip,}\\
     \texttt{\hspace*{2.4em}sp + n\_args);}\\
     \texttt{\hspace*{1.2em}ip = target;}\\
     \texttt{\hspace*{1.2em}DISPATCH\_JUMP();}\\
     \texttt{\}}};

\node[font=\large] (handlergap) at (8.1,-3.45) {$\cdots$};

\node[handler] (hconst) at (8.1,-4.35)
    {\texttt{void op\_const(DECL\_STATE) \{...\}}};

\node[handler] (hbegin) at (8.1,-5.50)
    {\texttt{void op\_begin(DECL\_STATE) \{...\}}};

\node[handler] (hend) at (8.1,-6.65)
    {\texttt{void op\_end(DECL\_STATE) \{...\}}};

\draw[flow] (ip.east) -- (callop.west);
\draw[flow] (callop.east) -- (hcall.west);
\draw[flow] (mainbegin.east) -- (hbegin.west);
\draw[flow] (constop.east) -- (hconst.west);
\draw[flow] (endop.east) -- (hend.west);

\draw[targetflow] (entrytarget1.east) -- (mainbegin.west);
\draw[targetflow] (entrygap.east) -- (gap2.west);

\draw[targetflow] (entrytargetn.east) -- (gap2.west);
\draw[targetflow] (calltarget.east) -- ++(0.55,0) |- (fbegin.east);
\draw[returnflow] (hend.west) to[out=180,in=0] (aftercall.east);

\end{tikzpicture}
}
    \caption{Исполнение внутреннего шитого кода виртуальной машины.}
    \label{fig:internal-threaded-code}
\end{figure}

Операции, создающие объекты Lama или обращающиеся к встроенным функциям,
используют существующую среду выполнения. Перед такими вызовами виртуальная
машина синхронизирует верхнюю границу своего стека со сборщиком мусора, чтобы
значения, лежащие на стеке, учитывались как корни.

Внешние вызовы реализованы через отдельную инструкцию вызова, использующую интерфейс внешних
функций. Имена внешних функций извлекаются из таблицы строк байткода и разрешаются через символы
текущего процесса. Для обычных функций используется единый механизм вызова. Отдельно обрабатываются
функции, которым среда выполнения передает массив аргументов, и вариативные функции наподобие
форматированного вывода. Такое разделение позволяет сохранить совместимость с существующими
соглашениями Lama, не добавляя специальные инструкции для каждой встроенной функции.

